% ============================================================================================================
% LipID v1.2 — experimental text skeleton. Every number is a \lipnum{key} generated by scripts/build_revision.py
% (results_v12/numbers_v12.tex, with definitions and sources in numbers_v12.csv). Load once in the preamble:
%   \input{results_v12/numbers_v12.tex}
% Wording in [brackets] is to be chosen AFTER the results exist, following the pre-specified reporting rules in
% protocols/v12_preregistration.json (small or null effects are reported as such).
% ============================================================================================================

% ------------------------------------------------------------------------------- Section 6.1 (unit of uncertainty)
\paragraph{Unit of uncertainty and the independent-seed batch.}
Within a seed, each $(d,\rho_c)$ shares one $\sigma$ draw, one set of 20 query/reference pairs and one nested anchor chain per
anchor law, so feature-instances are dependent and are never bootstrapped. Seeds 1--4 (exploration) and 8--11
(confirmation) keep their v1.1 role and are summarized descriptively. For this revision we locked a protocol
(\texttt{protocols/v12\_preregistration.json}, lock file and commit in the artifact) and then ran
\lipnum{stat-01.n_seeds} new independent seeds. Four pre-specified comparisons are estimated from one matched aggregate
per seed, with 95\% seed-level $t$ intervals, a seed bootstrap as sensitivity analysis and Holm adjustment within each
family (Table~\ref{tab:stat01}). These intervals describe variability over the random design of this synthetic study,
not uncertainty about real models or data.

% ------------------------------------------------------------------------------- Section 6.x (new) joint ranking
\subsection{Joint ranking information}\label{sec:jointexp}
Corollary~\ref{cor:pair} can only add certificates relative to comparing marginal intervals. We measure how many, on the
matched Section-6 instances (same seeds, anchor chains, query/reference pairs and Lipschitz levels; one row per feature
pair). At the certified level and $n=4096$ (data anchors, domain mode), the fraction of feature pairs with a strict signed
ordering rises from \lipnum{jnt-02.cert.d3.n4096.rate_marg} to \lipnum{jnt-02.cert.d3.n4096.rate_pair} at $d=3$ and from
\lipnum{jnt-02.cert.d6.n4096.rate_marg} to \lipnum{jnt-02.cert.d6.n4096.rate_pair} at $d=6$
(Table~\ref{tab:jnt02}, Fig.~\ref{fig:jnt02}); the seed-level gain is \lipnum{stat-01.P4.mean}
(95\% interval \lipnum{stat-01.P4.ci}). [Describe where the gain concentrates and where it vanishes; state explicitly
whether it is small near saturation.] Every true difference lies inside its sharp pairwise interval at every valid level
(\lipnum{jnt-02.integrity.not_contained_valid} exceptions), and no pairwise interval is wider than the marginal comparison
(\lipnum{jnt-02.integrity.v_dom} exceptions), as the corollary requires. Absolute rankings, which need the joint program
for $\phi_i+\phi_j$, are reported in the supplement (Table~\ref{tab:jnt03}). The closed form agrees with the
finite-hybrid linear program to the pre-specified $10^{-8}$ on every comparison of the validation grid
(Table~\ref{tab:jnt01}).

\begin{figure}[t]\centering
\includegraphics[width=\linewidth]{figures_v12/fig_jnt02_certification.pdf}
\caption{Strict signed pair certificates: sharp pairwise bounds (solid) versus marginal-interval comparison (dashed),
data anchors, domain mode. Lines are seed means; bands are 95\% seed-level $t$ intervals when at least ten seeds are
available, otherwise seed min--max (descriptive).}\label{fig:jnt02}
\end{figure}

% ------------------------------------------------------------------------------- Section 7.1 (MC, Q5)
\paragraph{Non-saturated instances.}
Because the Table-6 instances saturate for $m\ge12$, we repeated the comparison on instances whose anchors are placed by
a fixed, non-adaptive geometric rule near hybrid vertices (selected on a separate pilot, evaluated on independent
instances; Table~\ref{tab:scl02}). [Report excess width per method and budget, retained sign / ranking certificates and
cost; if every valid method remains near-vacuous, say so.] This design is a benchmark instance generator, not a proposed
query strategy.

% ------------------------------------------------------------------------------- Section 7.2 (metric, Q2)
\paragraph{Scale sensitivity.}
Intervals are invariant when all $\sigma_j$ and $L$ are rescaled together (verified to \lipnum{met-02.max_invariance}),
so only the \emph{relative} scales matter. With heterogeneous marginals and $L$ re-certified in each metric
(Table~\ref{tab:met02}), [describe the change in median $\rho$, saturation and certificates between MAD, SD and raw-unit
scales]. On the diabetes audit SD and raw units are close (the bundled features have equal SD on the full dataset, not on
the training split) (Table~\ref{tab:met02real}).

% ------------------------------------------------------------------------------- Section 7 (certificates, Q6)
\paragraph{Where the simple certificates lose saturation.}
Proposition~\ref{prop:chain} splits the gap between the exact normalized width and the geometry/range certificate into
direction asymmetry, slope magnitude and decoupling. Among saturated but uncertified features at $d=6$, $n=4096$
(domain mode, data anchors, $L^\star$), the shares first lost by each are \lipnum{cert-02.oracle.d6.sign_asymmetry},
\lipnum{cert-02.oracle.d6.slope_magnitude} and \lipnum{cert-02.oracle.d6.decoupling} (Table~\ref{tab:cert02}).
[Interpretation tied to the dominant class; no new certificate is proposed because the geometry-only pairwise bound is
more expensive than the exact envelope.]
